%!TEX program = uplatex
\RequirePackage{plautopatch}

\documentclass[uplatex,dvipdfmx]{jlreq}

% 余白調整（1ページ収め）
\usepackage[top=18mm,bottom=20mm,left=20mm,right=20mm]{geometry}

\usepackage{amsmath,amssymb}
\usepackage{url}

\pagestyle{empty}

% 行間少し詰める
\renewcommand{\baselinestretch}{0.96}

%--------------------------------
% 講演マクロ
%--------------------------------

\newcommand{\talk}[3]{%
\noindent
#1\quad #2\par
\hspace*{2em}#3\par\smallskip
}

%--------------------------------
% タイトル情報
%--------------------------------

\newcommand{\EwMnumber}{第2回}
\newcommand{\EwMtitle}{幾何学者は物理学から何を学んだか？}
\newcommand{\EwMdate}{1997年2月21日（金）14:30 ～ 2月22日（土）16:30}
\newcommand{\EwMplace}{東京都文京区春日1-13-27 中央大学理工学部5号館 5534号室}

%--------------------------------
% タイトル表示
%--------------------------------

\newcommand{\EwMtitleblock}{

\vspace*{-5mm}

\begin{center}

{\Large ENCOUNTER with MATHEMATICS}

\vspace{2mm}

{\large 数学との遭遇}

\vspace{4mm}

{\Large \bfseries \EwMnumber\ \EwMtitle}

\vspace{4mm}

\EwMdate

\vspace{2mm}

於：\EwMplace

\end{center}

\vspace{3mm}

\vspace{2mm}

}

\begin{document}

\EwMtitleblock

%--------------------------------
% プログラム
%--------------------------------

\section*{プログラム}

\subsection*{1日目（2月21日（金））}

\talk
{14:30～15:40}
{数学と物理のこの20年 I}
{深谷 賢治 氏（京大・理）}

\talk
{16:20～17:30}
{Donaldson 不変量の構造定理への4つのアプローチ I}
{古田 幹雄 氏（京大・数理研）}

\subsection*{2日目（2月22日（土））}

\talk
{10:30～11:45}
{数学と物理のこの20年 II}
{深谷 賢治 氏（京大・理）}

\talk
{13:20～14:35}
{Donaldson 不変量の構造定理への4つのアプローチ II}
{古田 幹雄 氏（京大・数理研）}

\talk
{15:10～16:25}
{数学と物理のこの20年 III}
{深谷 賢治 氏（京大・理）}

%--------------------------------
% 案内文
%--------------------------------

\bigskip

別紙の趣旨に沿った集会の第2回を以上のような予定で開催いたします。概ね非専門家向けということで講演をお願い致しましたので、多くの方々のご参加をお待ちしております。講演者による講演内容へのご案内を別紙に添付いたしますのでご覧下さい。

%--------------------------------
% 連絡先
%--------------------------------

\section*{連絡先}

112 東京都文京区春日1-13-27 中央大学理工学部数学教室

tel : 03-3817-1745

ENCOUNTER with MATHEMATICS  
e-mail : encmath@math.chuo-u.ac.jp  

三松 佳彦 : e-mail : yoshi@math.chuo-u.ac.jp  

\url{http://www.math.chuo-u.ac.jp/ENCwMATH}

\end{document}
